\section{Open Problems and Conclusion}\label{sec:outlook}

The loop in \cref{fig:loop} is now closed at every layer, but each layer closes it against its own
metric: the radio against delay budgets, transport against frame loss, the renderer against image
quality. Five questions would join the parts.

\begin{openproblems}
\item \textbf{Do QoE thresholds survive real network dynamics?} Every user study in \cref{sec:qoe} applied
  static delay, independent loss or fixed jitter, while measured wireless loss arrives in bursts
  \cite{W7202080971}. Replaying recorded Wi-Fi and 5G traces to participants would show whether the
  tolerance of 75--77~ms RTT and the dominance of loss still hold.
\item \textbf{Does frame-aware scheduling raise QoE?} Frame-aware schedulers report 15--45\% more capacity
  \cite{W4402156740,W4406264804}, counted as packets that meet a PDB under traffic models that
  underestimate real demand \cite{W4399340506}. A testbed that couples such a scheduler to a headset and a
  QoE model would test whether capacity gains are experience gains.
\item \textbf{Who should own the rate?} Network recommendations \cite{W7161282618}, Release-19 uplink rate
  control \cite{W7164211432}, L4S marks \cite{W4408870658} and application bitrate adaptation
  \cite{W4405014964} can all act on one flow. Their joint stability, fairness and response time are
  unstudied.
\item \textbf{What is group QoE?} Social studies test pairs \cite{W4392600587,W4415821547}, while
  platforms serve five or more users through forwarding servers \cite{W4403981191}. How delay assigned to
  one member changes the experience of the others, and how to share impairments fairly, is unknown.
\item \textbf{Which representation for which link?} Point clouds, keypoint avatars, Gaussian splats and
  remotely rendered video trade bandwidth, device computation and latency differently, but each is
  evaluated alone and on image metrics \cite{W4404525293,W4416477288,W7160906632}. A common user-facing
  benchmark over the same network traces is missing.
\end{openproblems}

Since 2023, networked XR has gained frame-aware radios, earlier congestion signals, transports with
deadlines and content built for adaptation, while users have shown that reliability matters more than the
last tens of milliseconds. The next step is to judge each mechanism by the experience it produces under
real network dynamics, alone and in groups.
